% !TeX root = ../../Book4.tex
% 纲要标题：### 19.3 心
% 纲要位置：计划纲要.md，第 2993 行。

\section{心}
\label{b4:sec:1903}


% 纲要内容（仅作写作计划，尚非教材正文）：
%
% * 心是交换范畴
% * 短正合列与三角的关系
% * 标准心与原交换范畴
%

\subsection{心是交换范畴}
\label{b4:subsec:1903:01}

在心中，一切对象都位于同一个上同调次数；但两个对象之间态射的锥通常占据相邻两次。这正好提供核与余核所需的两个对象。困难在于还要证明余像等于像，不能只造出两个候选对象便结束。

\begin{theorem}\label{b4:thm:heart-abelian}
设 \(\mathcal D\) 为带 t-结构的三角范畴。其心 \(\mathcal H\) 是交换范畴。若 \(f:X\to Y\) 是心中态射，完成为好三角
\[
 X\xrightarrow{f}Y\longrightarrow C\longrightarrow X[1],
\]
则心中的核、余核分别为
\[
 \ker_{\mathcal H}f=H_t^{-1}(C),\qquad
 \operatorname{coker}_{\mathcal H}f=H_t^0(C).
\]
这两个对象及其核、余核映射由泛性质唯一确定，不依赖锥的选择。
\end{theorem}
\begin{proof}
心是满的加性子范畴，因为\cref{b4:prop:t-orthogonals-extensions}保证有限直和仍在心中。旋转三角 \(Y\to C\to X[1]\) 表明 \(C\in\mathcal D^{\ge-1}\cap\mathcal D^{\le0}\)。令
\[
 K=(\tau^{\le-1}C)[-1],\qquad Q=\tau^{\ge0}C.
\]
两者都在心中，有截断三角 \(K[1]\to C\to Q\to K[2]\)。其与原三角给出 \(K\to X\) 及 \(Y\to Q\)。

对心中任意 \(T\)，在原三角上取 \(\operatorname{Hom}(T,-)\)，得到
\[
 0\longrightarrow\operatorname{Hom}(T,C[-1])
 \longrightarrow\operatorname{Hom}(T,X)
 \longrightarrow\operatorname{Hom}(T,Y),
\]
首项前的 \(\operatorname{Hom}(T,Y[-1])\) 为零。截断三角及 \(Q[-1],Q[-2]\) 的正次数条件又给出 \(\operatorname{Hom}(T,K)\simeq\operatorname{Hom}(T,C[-1])\)。所以 \(K\to X\) 是核。对偶地，\(\operatorname{Hom}(Q,T)\simeq\operatorname{Hom}(C,T)\)，而 \(\operatorname{Hom}(X[1],T)=0\)，因此 \(Y\to Q\) 是余核。

还须核对余像与像。对复合 \(Y\to C\to Q\) 应用八面体公理，取该复合的锥为 \(I[1]\)，得到相容的两个好三角
\[
 K\longrightarrow X\longrightarrow I\longrightarrow K[1],
 \qquad
 I\longrightarrow Y\longrightarrow Q\longrightarrow I[1],
\]
其中 \(X\to I\to Y\) 的复合就是 \(f\)。第一三角使 \(I\in\mathcal D^{\le0}\)，第二三角使 \(I\in\mathcal D^{\ge0}\)，故 \(I\) 在心中。在第一三角上向任意心对象取 Hom，\(\operatorname{Hom}(K[1],T)=0\) 表明 \(X\to I\) 是 \(K\to X\) 的余核。在第二三角上从心对象取 Hom，\(\operatorname{Hom}(T,Q[-1])=0\) 表明 \(I\to Y\) 是 \(Y\to Q\) 的核。于是 \(I\) 同时给出余像与像，且它们间的典范态射是同构。这验证了交换范畴的全部条件。
\end{proof}

上述八面体只用于比较两个实际构造出的三角：一边先去掉 \(f\) 的核，另一边先去掉它的余核。它所产生的共同对象 \(I\) 才使两种做法一致。原始证明可对照\cite[Proposition 1.2.4; Théorème 1.3.6]{BBD1982FaisceauxPervers}。

\begin{theorem}\label{b4:thm:t-cohomology-exact}
设 \(\mathcal D\) 带 t-结构，心为 \(\mathcal H\)。函子 \(H_t^0:\mathcal D\to\mathcal H\) 是上同调函子：每个好三角 \(X\to Y\to Z\xrightarrow{w}X[1]\) 诱导自然长正合列
\[
 \cdots\longrightarrow H_t^n(X)\longrightarrow H_t^n(Y)
 \longrightarrow H_t^n(Z)\xrightarrow{H_t^n(w)}H_t^{n+1}(X)
 \longrightarrow\cdots.
\]
\end{theorem}
\begin{proof}
只须证明 \(H_t^0(X)\to H_t^0(Y)\to H_t^0(Z)\) 正合，再平移并旋转三角。

先设三项都在 \(\mathcal D^{\le0}\)。对 \(T\in\mathcal H\)，伴随给出 \(\operatorname{Hom}(U,T)=\operatorname{Hom}(H_t^0(U),T)\)，其中 \(U=X,Y,Z\)。Hom 长正合列因 \(\operatorname{Hom}(X[1],T)=0\) 而给出
\[
 0\longrightarrow\operatorname{Hom}(H_t^0Z,T)
 \longrightarrow\operatorname{Hom}(H_t^0Y,T)
 \longrightarrow\operatorname{Hom}(H_t^0X,T).
\]
这正是余核泛性质，所以 \(H_t^0X\to H_t^0Y\to H_t^0Z\to0\) 正合。

再设仅有 \(X\in\mathcal D^{\le0}\)。对 \(T\in\mathcal D^{\ge1}\)，正交性使 \(\operatorname{Hom}(Z,T)\to\operatorname{Hom}(Y,T)\) 同构。因此 \(\tau^{\ge1}Y\to\tau^{\ge1}Z\) 是同构。把 \(Y\to Z\to\tau^{\ge1}Z\) 代入八面体公理，其余两锥分别是截断的低次部分的移位，得到好三角
\[
 X\longrightarrow\tau^{\le0}Y\longrightarrow\tau^{\le0}Z\longrightarrow X[1].
\]
应用前一情形得到 \(H_t^0X\to H_t^0Y\to H_t^0Z\to0\) 正合。对偶的论证表明：若 \(Z\in\mathcal D^{\ge0}\)，则 \(0\to H_t^0X\to H_t^0Y\to H_t^0Z\) 正合。这里对偶把三角范畴换为反范畴，交换两类次数条件，故完全使用同样的伴随与 Hom 正合性。

一般情形下，令 \(A=\tau^{\le0}X\)、\(B=\tau^{\ge1}X\)。对 \(A\to X\to Y\) 应用八面体，得到三角 \(A\to Y\to U\) 与 \(U\to Z\to B[1]\)。第一三角给出 \(H_t^0X\to H_t^0Y\to H_t^0U\to0\) 正合；因 \(B[1]\in\mathcal D^{\ge0}\)，第二三角给出单态射 \(H_t^0U\to H_t^0Z\)。复合后正是所需的三项正合性。连接映射及其自然性来自原三角的态射，因而各次正合列可接成所述长列。
\end{proof}

\subsection{短正合列与三角的关系}
\label{b4:subsec:1903:02}

\begin{proposition}\label{b4:prop:heart-exact-triangles}
设 \(\mathcal H\) 为三角范畴 \(\mathcal D\) 中一个 t-结构的心。心中的列 \(0\to A\xrightarrow{i}E\xrightarrow{p}B\to0\) 短正合，当且仅当存在 \(\delta:B\to A[1]\) 使 \(A\xrightarrow{i}E\xrightarrow{p}B\xrightarrow{\delta}A[1]\) 为好三角。
\end{proposition}
\begin{proof}
若三角存在，\cref{b4:thm:t-cohomology-exact}的长正合列因三项都在心中而成为所述短列。反过来，将 \(i\) 完成为三角 \(A\to E\to C\)。\cref{b4:thm:heart-abelian}给出 \(H_t^{-1}(C)=\ker i=0\) 和 \(H_t^0(C)=\operatorname{coker}i=B\)。又 \(C\in\mathcal D^{[-1,0]}\)，故其截断三角直接给出 \(C\simeq B\)，并把原来的余核映射识别为 \(p\)。
\end{proof}

\begin{proposition}\label{b4:prop:heart-ext1}
设 \(A,B\) 属于 t-结构的心 \(\mathcal H\)。心中的一次 Yoneda 扩张与环境中的移位态射自然对应：
\[
 \operatorname{YExt}_{\mathcal H}^1(B,A)
       \simeq\operatorname{Hom}_{\mathcal D}(B,A[1]).
\]
左边采用短正合列的等价类与 Baer 和；若心有足够内射或投射对象，它就是通常的一次 Ext。
\end{proposition}
\begin{proof}
短正合列对应前一命题中的连接映射。它由给定短列唯一确定：若有两个三角结构，三角态射公理与心中的核、余核泛性质使三个零次对象上的比较映射都是恒等。反之，给定 \(\delta:B\to A[1]\)，完成并旋转三角得到 \(A\to E\to B\xrightarrow{\delta}A[1]\)。心对扩张封闭，故 \(E\in\mathcal H\)，前一命题给出短正合列。具有同一 \(\delta\) 的两个三角之间可用三角态射公理作比较，中间映射为同构，因而所得扩张等价。两种操作互逆。

对端点的态射取推出或拉回，三角态射的最后方块表明连接映射相应地作后复合或前复合。直和扩张对应连接映射的直和。Baer 和先沿 \(B\to B\oplus B\) 拉回，再沿 \(A\oplus A\to A\) 推出，因此对应于两个连接映射相加。这也证明群结构及对两端的自然性。
\end{proof}

这里把扩张送到三角的连接箭头 \(\delta\)，与\secref{b4:sec:1704}经内射消解识别一次 Ext 的约定一致。用其他符号约定书写锥时，仍须同时核对连接箭头与消解中的余圈代表。

\ProseParagraphBreak
高阶扩张可以把一次扩张拼接成链，再把连接态射移位后复合，从而得到到 \(\operatorname{Hom}_{\mathcal D}(B,A[n])\) 的比较映射；它在扩张态射下不变，因为每个方块的连接态射自然。一次的双射不能直接推广为每个 \(n\) 的双射。\secref{b4:sec:1905}将给出二次比较失效的具体例子。

\subsection{标准心与原交换范畴}
\label{b4:subsec:1903:03}

\begin{proposition}\label{b4:prop:standard-heart-equivalence}
设 \(\mathcal A\) 为交换范畴。把对象视为集中于零次的复形，得到 \(\mathcal A\) 到 \(D(\mathcal A)\) 标准心的正合等价。
\end{proposition}
\begin{proof}
若 \(X\) 只有 \(H^0\) 可能非零，截断给出拟同构屋顶 \(X\leftarrow\tau^{\le0}X\to H^0(X)[0]\)，所以该函子本质满。对 \(A,B\in\mathcal A\)，用屋顶 \(A\leftarrow Z\to B\) 表示导出态射；先以 \(\tau^{\le0}Z\) 细化，它到 \(A,B\) 的映射都经 \(H^0(Z)[0]\) 分解，且 \(H^0(Z)\to A\) 是同构。因此态射唯一来自 \(A\to B\)。唯一性也可直接对该屋顶取 \(H^0\) 看出。这证明全忠实。最后，原范畴的短正合列给出好三角，再由\cref{b4:prop:heart-exact-triangles}成为心中短正合列，所以等价正合。
\end{proof}

\begin{proposition}\label{b4:prop:nondegenerate-cohomology-conservative}
设 \(\mathcal D\) 的 t-结构非退化。若 \(H_t^n(X)=0\) 对所有整数 \(n\) 成立，则 \(X=0\)；若态射 \(f\) 在所有 \(H_t^n\) 上诱导同构，则 \(f\) 是同构。
\end{proposition}
\begin{proof}
先设 \(X\in\mathcal D^{\le0}\)。相邻截断三角和各次上同调的消失依次使 \(X\in\mathcal D^{\le-1},\mathcal D^{\le-2},\ldots\)，故非退化性给出 \(X=0\)。对 \(X\in\mathcal D^{\ge1}\) 对偶成立。一般 \(X\) 的 \(\tau^{\le0}X\) 与 \(\tau^{\ge1}X\) 仍有所有上同调为零，分别用前两种情形即可。若 \(f\) 诱导全部上同调同构，其锥由长正合列知全部上同调为零，因而锥为零，\(f\) 为同构。
\end{proof}

“检测同构”不等于“用所有上同调恢复每个态射”。例如非零的 \(\operatorname{Ext}_{\mathbb Z}^1(\mathbb Z/2,\mathbb Z)\) 给出非零导出态射 \(\mathbb Z/2\to\mathbb Z[1]\)，但它在每一个相同次数的普通上同调之间都诱导零映射。
